\begin{abstract}
Rich annotation schemes specify structures and conventions that a general-purpose LLM may not follow. We study whether supervised revision of a modular annotation harness can improve annotation from a few gold documents without changing model weights. \sys records step traces, localizes discrepancies using structured gold alignment, and applies typed, anchored edits to the instructions of designated skills. Proposed tools and steps undergo machine checks and a separate model review. In one fixed-order Chinese UMR run, a paired test-then-train evaluation gives a mean gain of \gainF{} sentence-graph F-score points over the seed on \nPaired{} previously untrained-on documents (document bootstrap 95\% interval \gainFLo--\gainFHi). Gains concentrate in convention-sensitive checks; little change is observed in sense, aspect and core-argument checks. Stage-wise scores associate most of the gain with the generator, while the checking stage adds \critFirst{} points. No tool proposal is admitted. Instructions grow \sizeFactorInt-fold without a corresponding increase in sequential gains, and a separate 9B joint builder--annotator configuration degrades. These are observations from one trajectory, rather than evidence of causal attribution or broad superiority. Matched held-out baselines, design ablations and repeated runs remain pending.
\end{abstract}
